MHC-I presentation can vary among tissues even when the presenting allele and source protein are fixed. Existing predictors largely estimate tissue-blind peptide--MHC compatibility and therefore do not directly address this tissue-associated preference. We formulate a matched-pair ranking problem: given two peptides from the same parent protein under the same presenting restriction, determine which has stronger presentation evidence in a specified tissue. To remove a simple cross-tissue-frequency cue, the two peptides in every pair are additionally matched on their total number of positive tissue occurrences. The resulting benchmark contains 77 Human tissue--HLA tasks and 11 Mouse tissue--H-2 tasks, where H-2 is the mouse major histocompatibility complex. Both species use one fixed internal train/test split with 50 test pairs per task and the same three final-training seeds (20260704--20260706). Species-specific TissuePMHC configurations are selected by three-fold cross-validation confined to the training files before the fixed tests are opened. The occurrence audit finds zero mismatched pairs and identical positive-occurrence distributions between labels in both species and partitions. Tuned TissuePMHC reaches mean AUROC/AUPRC values of 0.8136/0.8126 on Human and 0.8194/0.8279 on Mouse. Same-condition ablations support the multi-kernel encoder in both species, while a trainable MHC-only control is worse than TissuePMHC in every task. Post-hoc Expected-integrated-gradients SHAP analysis of the Human model localizes the largest pair-matched sequence contributions to peptide positions 2, 3, and 9 and reveals stable HLA-stratified residue patterns. Exhaustive single-residue in silico mutagenesis independently recovers the same position ordering and agrees with observed-residue SHAP attribution in the global and HLA-specific branches (three-seed ensemble Spearman \(\rho=0.827\) and \(0.911\)). When the peptide, HLA, and substitution are fixed and only the tissue task is changed, the median HLA-level tissue-to-seed perturbation-dispersion ratio is 1.899, with tissue dispersion exceeding seed dispersion for all 28 multi-tissue HLA alleles. Tissue-independent external predictors remain near random in Human (AUROC 0.4874--0.5044) and Mouse (0.4917--0.5009). The tuning folds are not treated as strict-generalization estimates, and no peptide-separated evaluation is reported. These results assess tissue-conditioned peptide ranking for previously represented tissue--MHC tasks; the interpretation analyses describe model dependence and do not establish a causal tissue-processing mechanism, unseen-context generalization, independent-cohort validity, or clinical utility.
